\documentclass[12pt]{article}
\usepackage[utf8]{inputenc}
\usepackage[T1]{fontenc}
\usepackage[margin=1in]{geometry}
\usepackage{setspace}
\usepackage{xcolor}
\usepackage[colorlinks=true, urlcolor=blue!60!black, linkcolor=blue!60!black]{hyperref}
\setlength{\parindent}{0em}
\setlength{\parskip}{0.6em}
\renewcommand{\baselinestretch}{1.0}
\begin{document}
\begin{center}
{\LARGE \textit{Jacobson v. Massachusetts}, 197 U.S. 11 (1905)\par}
\vspace{0.5em}
{\large Case Brief --- Two-Page Edition\par}
\end{center}

\section{Memory Jogger}
States may compel adult vaccination during an epidemic: liberty is subject to reasonable restraints for the common good (\href{https://supreme.justia.com/cases/federal/us/197/11/}{Justia}).

\section{Detailed Case Facts}
During a 1902 smallpox outbreak, the Cambridge Board of Health, acting under a state statute, ordered all inhabitants vaccinated, with a \$5 fine for adults who refused. Pastor Henning Jacobson refused, citing past adverse reactions, but offered no proof he was then medically unfit.

\section{Procedural History}
The trial court excluded his fourteen offers of proof against vaccination; he was convicted and fined \$5, and the Supreme Judicial Court affirmed.

\section{Judicial Votes}
Affirmed 7--2. Harlan wrote for the Court; Brewer and Peckham dissented without opinion.

\section{Holding}
A State may, under its police power, require vaccination during an outbreak. Courts may strike down a health law only if it has ``no real or substantial relation'' to its object or is a ``plain, palpable invasion of rights secured by the fundamental law.''

\section{Analysis of Opinions}
\textbf{Majority:} Liberty is not freedom from all restraint; quarantine and conscription show the community's power of self-defense. The legislature, not a jury, chooses among medical theories. Children and adults could be treated differently. But a law ``arbitrary and oppressive'' as applied---such as forcing vaccination on someone it would likely kill---would warrant judicial intervention. \textbf{Dissents:} Brewer and Peckham wrote nothing, but joined forces months later in \textit{Lochner}.

\section{Examples: Future Applications}
\textbf{Same-Side:} school immunization mandates (\textit{Zucht v. King}). \textbf{Opposite-Side:} capping worship services while sparing secular venues (\textit{Roman Catholic Diocese v. Cuomo}). \textbf{Fence-Sitter:} indefinite heavy fines for declining boosters against a low-fatality endemic virus.

\section{Critique}
Critics say courts turned a modest \$5-fine decision into an emergency ``super-precedent,'' and note that \textit{Buck v. Bell} relied on it. Public-health scholars read it as a balanced framework of necessity and proportionality.

\section{Key Quotations}
\begin{quote}
``But the liberty secured by the Constitution of the United States to every person within its jurisdiction does not import an absolute right in each person to be, at all times and in all circumstances, wholly freed from restraint. There are manifold restraints to which every person is necessarily subject for the common good.''
\end{quote}
\noindent --- Justice Harlan, for the Court, 197 U.S. at 26.

\end{document}
